% P15-U1: extended means, renewal, empirical distributions and finite information.
\chapter{Applications of Strong Laws}
\label{sec:p15-strong-law-applications-chapter}

Strong laws also control inverse counting processes, an entire empirical
distribution function and the probability of a long information string.
Each application requires its own conversion from an average to the
claimed object. The treatment below supplies those conversions and their
measurability and endpoint qualifications.

The source is Durrett's author Version 5 dated 11 January 2019,
Theorem 2.4.5, the renewal construction and Theorem 2.4.7, the
random-index observation, Theorem 2.4.9, and the finite-alphabet
information Example 2.4.10 \cite[pp.~78--80]{DurrettPTEAuthor2019}.
Zero waiting times and zero-probability symbols are explicit extensions
of the source's strictly positive regimes; their additional arguments
are proved below. The quantitative typical-set bounds are also derived
here.

\section{Extended means and inverse counts}
\label{sec:p15-extended-means-renewal}

\begin{hbdefinition}{def:p15-extended-mean-law}{The domain of an extended expectation}
On a probability space \((\Omega,\mathcal F,\mathbb P)\), for a
real random variable \(X\), let
\(X^+=\max(X,0)\) and \(X^-=\max(-X,0)\). Its expectation
\(\mathbb EX=\mathbb EX^+-\mathbb EX^-\) is defined as an extended
real number when at least one of these two nonnegative expectations is
finite. If both are infinite, no expectation is assigned by cancellation.
For a sequence \(X_n\), write \(S_n=\sum_{k=1}^nX_k\). A limit
\(S_n/n\to+\infty\) means that for each finite real \(K\), all
sufficiently late averages exceed \(K\); the negative limit is defined
by negation.
\end{hbdefinition}

\begin{hbtheorem}{thm:p15-extended-mean-strong-law}{proved=annex:p15.extended-mean-strong-law}{The strong law whenever the extended mean exists}
Let \(X_1,X_2,\ldots\) be i.i.d. real random variables.
If their expectation exists in the preceding sense, then
\[
 \frac{S_n}{n}\longrightarrow\mathbb EX_1
                                  \qquad\text{almost surely}.
\]
In particular, \(\mathbb EX_1^+=\infty\) and
\(\mathbb EX_1^-<\infty\) give limit \(+\infty\), while the reversed
conditions give limit \(-\infty\). When both part expectations are
finite, the conclusion is the integrable strong law. The assertion
does not cover the case in which both are infinite.
\end{hbtheorem}

\begin{hblemma}{thm:p15-random-index-limits}{proved=annex:p15.random-index-limits}{Almost-sure limits under random indices}
If measurable extended real \(Z_n\to Z\) almost surely and finite
positive-integer-valued random variables \(J_m\to\infty\) almost
surely, then \(Z_{J_m}\) is measurable and \(Z_{J_m}\to Z\) almost
surely. Independence between the indices and the observations is not
required.

The corresponding assertion for convergence in probability fails.
For independent Bernoulli variables \(B_n\) with
\(\mathbb P(B_n=1)=a_n\), where \(0\leq a_n\leq1\),
\(a_n\to0\) and \(\sum_na_n=\infty\), we have \(B_n\to0\)
in probability. One choice is \(a_n=1/(n+1)\).
Let \(J_m\) be the first \(j\geq m\) with \(B_j=1\), assigning
\(J_m=m\) if no such \(j\) exists. These indices are finite measurable
variables with \(J_m\geq m\), but \(B_{J_m}=1\) for every \(m\)
on one event of probability one.
\end{hblemma}

\begin{hbdefinition}{def:p15-renewal-count}{Renewal epochs, ties and the inverse count}
Let \(X_1,X_2,\ldots\) be i.i.d. nonnegative finite real waiting
times. Set
\[
 T_0=0,\qquad T_n=\sum_{k=1}^nX_k,\qquad
 N(t)=\sup\{n\in\mathbb N_0:T_n\leq t\}\quad(t\geq0),
\]
where \(\mathbb N_0=\{0,1,\ldots\}\) and an unbounded set of indices
has supremum \(+\infty\). A zero waiting time permits coincident
epochs. The definition counts every renewal at or before \(t\), so
it chooses the last index in a finite group of tied epochs.
\index{renewal process}\index{renewal count}
\end{hbdefinition}

\begin{hbtheorem}{thm:p15-renewal-inverse-law}{proved=annex:p15.renewal-inverse-law}{The renewal count with zero waiting times allowed}
Assume \(\mathbb P(X_1>0)>0\) in the preceding definition, and put
\(\mu=\mathbb EX_1\). Then \(0<\mu\leq\infty\).
For each fixed \(t\), \(N(t)\) is an extended integer-valued
measurable variable. There is one measurable event of probability one
on which \(T_n\to\infty\), every \(N(t)\) is finite,
\(N(t)\to\infty\) as real \(t\to\infty\), and
\[
 \frac{N(t)}t\longrightarrow\frac1\mu,
 \qquad 1/\infty:=0.
\]
On that event,
\(T_{N(t)}\leq t<T_{N(t)+1}\), including when epochs coincide.
Division by \(N(t)\) is used only for sufficiently large \(t\),
when \(N(t)\geq1\).

If instead \(\mathbb P(X_1>0)=0\), then \(T_n=0\) for all \(n\)
and \(N(t)=\infty\) for every \(t\geq0\) almost surely. This
degenerate case is not a finite renewal-count process.
\end{hbtheorem}

\section{Uniform empirical distributions}
\label{sec:p15-empirical-distributions}

\begin{hbdefinition}{def:p15-empirical-cdf}{Empirical measures, both threshold conventions and the uniform error}
Let \(X_n\) be i.i.d. real observations with common law \(\nu\)
and cumulative distribution function \(F(x)=\nu((-\infty,x])\).
For each \(n\geq1\) and outcome, define the empirical probability
measure and its distribution function by
\[
 \nu_n(B)=\frac1n\sum_{i=1}^n\mathbf1_{\{X_i\in B\}},
 \qquad
 F_n(x)=\nu_n((-\infty,x]).
\]
Here \(B\) ranges over Borel subsets of \(\mathbb R\).
Use the left-limit notation
\[
 F(x-)=\nu((-\infty,x)),\qquad
 F_n(x-)=\frac1n\sum_{i=1}^n\mathbf1_{\{X_i<x\}}.
\]
At formal grid endpoints, set both function values and left-limit
values at \(-\infty\) to zero, and those at \(+\infty\) to one.
The uniform CDF error is
\[
 D_n=\sup_{x\in\mathbb R}|F_n(x)-F(x)|.
\]
For an integer \(k\geq2\), the quantile grid is
\[
 x_{j,k}=\inf\{y\in\mathbb R:F(y)\geq j/k\}
 \quad(1\leq j<k),\qquad x_{0,k}=-\infty,\quad x_{k,k}=+\infty.
\]
Repeated interior grid points are allowed.
\index{empirical distribution function}\index{Glivenko--Cantelli theorem}
\end{hbdefinition}

\begin{hblemma}{thm:p15-empirical-cdf-structure}{proved=annex:p15.empirical-cdf-structure}{Measurable uniform error and the quantile inequalities}
For each outcome \(\nu_n\) is a probability measure and \(F_n\)
is its right-continuous CDF, with the stated left limits.
The random variable \(D_n\) is measurable, takes values in \([0,1]\),
and satisfies
\[
 D_n=\sup_{q\in\mathbb Q}|F_n(q)-F(q)|.
\]
Each interior \(x_{j,k}\) is finite, the grid is nondecreasing, and
\[
 F(x_{j,k}-)\leq j/k\leq F(x_{j,k}).
\]
For every adjacent pair, including the infinite endpoints,
\[
 F(x_{j,k}-)-F(x_{j-1,k})\leq1/k\quad(1\leq j\leq k).
\]
\end{hblemma}

\begin{hblemma}{thm:p15-continuous-cdf-uniform}{proved=annex:p15.continuous-cdf-uniform}{The deterministic continuous-CDF case}
Let \(G_n\) be distribution functions and let \(F\) be a continuous
distribution function on \(\mathbb R\). If \(G_n(x)\to F(x)\)
at every real \(x\), then
\(\sup_x|G_n(x)-F(x)|\to0\).
The condition that the limiting function is a distribution function,
including its two tail limits, is part of this statement.
\end{hblemma}

\begin{hbtheorem}{thm:p15-glivenko-cantelli}{proved=annex:p15.glivenko-cantelli}{The Glivenko--Cantelli theorem for arbitrary real laws}
For the i.i.d. observations in
Definition~\ref{def:p15-empirical-cdf},
\[
 \sup_{x\in\mathbb R}|F_n(x)-F(x)|\longrightarrow0
                                  \qquad\text{almost surely}.
\]
No continuity or density assumption on \(F\) is required. On one
event of probability one, for every integer \(k\geq2\) there is an
index \(n_k\) such that \(D_n\leq2/k\) whenever \(n\geq n_k\).
The index is obtained from convergence at the finitely many grid points
and their left limits; the proof includes grid points themselves and
empty intervals caused by repetitions.
\end{hbtheorem}

\section{Finite-alphabet information and typical sets}
\label{sec:p15-finite-alphabet-information}

\begin{hbdefinition}{def:p15-finite-alphabet-source}{Word probabilities, entropy and zero-mass symbols}
Let \(\mathcal A\) be a nonempty finite alphabet with its full
power-set sigma-algebra, let
\(p(a)\geq0\) with \(\sum_{a\in\mathcal A}p(a)=1\), and let
\(X_n\) be independent with this common law.
Set \(\mathcal A_+=\{a:p(a)>0\}\) and \(m=|\mathcal A_+|\geq1\).
For a word \(w=(w_1,\ldots,w_n)\in\mathcal A^n\), put
\[
 \pi_n(w)=\prod_{i=1}^np(w_i),\qquad W_n=(X_1,\ldots,X_n).
\]
All logarithms in this section are natural logarithms; information is
therefore measured in nats. Set \(-\log0=+\infty\) for a word of
zero probability, and define entropy by
\[
 H=-\sum_{a\in\mathcal A_+}p(a)\log p(a).
\]
Equivalently the sum may include zero-mass symbols with the convention
\(0\log0=0\).
For an everywhere real measurable observation, define
\[
 I_0(a)=
 \begin{cases}-\log p(a),&a\in\mathcal A_+,\\0,&a\notin\mathcal A_+,\end{cases}
 \qquad
 \sigma_I^2=\sum_{a\in\mathcal A_+}p(a)(-\log p(a)-H)^2.
\]
The value of \(I_0\) at a zero-mass symbol is a null-set convention,
not the self-information of an impossible symbol.
\index{entropy!finite alphabet}\index{asymptotic equipartition}
\end{hbdefinition}

\begin{hblemma}{thm:p15-finite-alphabet-entropy}{proved=annex:p15.finite-alphabet-entropy}{Normalization, entropy bounds and the information sum}
The function \(\pi_n\) is the probability mass function of \(W_n\).
On one event of probability one all \(X_i\) belong to
\(\mathcal A_+\), and for every \(n\),
\[
 -\log\pi_n(W_n)=\sum_{i=1}^n I_0(X_i).
\]
The observations \(I_0(X_i)\) are i.i.d. bounded real variables with
mean \(H\) and variance \(\sigma_I^2<\infty\). Moreover,
\[
 0\leq H\leq\log m.
\]
In particular, \(m=1\) gives \(H=\sigma_I^2=0\).
\end{hblemma}

\begin{hbtheorem}{thm:p15-finite-alphabet-aep}{proved=annex:p15.finite-alphabet-aep}{The finite-alphabet asymptotic equipartition property}
For the specified independent source,
\[
 -\frac1n\log\pi_n(W_n)\longrightarrow H
                                  \qquad\text{almost surely}.
\]
For every real \(\varepsilon>0\), define its typical set
\[
 \mathcal T_n(\varepsilon)=
 \{w\in\mathcal A^n:
 e^{-n(H+\varepsilon)}\leq\pi_n(w)
                  \leq e^{-n(H-\varepsilon)}\}.
\]
Then \(\mathbb P(W_n\in\mathcal T_n(\varepsilon))\to1\).
Zero-probability words never belong to this set, and the assertion
includes a one-symbol support. The inequalities defining typicality
are non-strict.
\end{hbtheorem}

\begin{hbtheorem}{thm:p15-typical-set-bounds}{proved=annex:p15.typical-set-bounds}{Explicit probability and cardinality bounds}
Put \(\eta_n=\mathbb P(W_n\notin\mathcal T_n(\varepsilon))\).
For every \(n\geq1\) and \(\varepsilon>0\),
\[
 \eta_n\leq\min\left\{1,\frac{\sigma_I^2}{n\varepsilon^2}\right\},
 \qquad
 (1-\eta_n)e^{n(H-\varepsilon)}
       \leq|\mathcal T_n(\varepsilon)|\leq e^{n(H+\varepsilon)}.
\]
For every subset \(\mathcal C\subseteq\mathcal A^n\),
\[
 \mathbb P(W_n\in\mathcal C)
       \leq\eta_n+|\mathcal C|e^{-n(H-\varepsilon)}.
\]
Thus if \(\mathbb P(W_n\in\mathcal C)\geq1-\delta\), where
\(0<\delta<1\), then
\[
 |\mathcal C|\geq\max\{0,1-\delta-\eta_n\}\,e^{n(H-\varepsilon)}.
\]
In particular \(n\geq\sigma_I^2/(\delta\varepsilon^2)\), with
\(n\geq1\), ensures typical-set probability at least \(1-\delta\).
These finite bounds are derived consequences of the information-sum
representation, beyond the source's stated limiting probability.
\end{hbtheorem}

\section{Computed applications and failure cases}
\label{sec:p15-strong-law-application-examples}

\begin{hbtheorem}{thm:p15-strong-law-application-calculations}{proved=annex:p15.strong-law-application-calculations}{Renewal ties, an atomic CDF and a biased alphabet}
Use i.i.d. waiting-time and observation sequences in the following
examples, and the independent alphabet source defined above. Their full
calculations are in the annex.
\begin{enumerate}
\item Waiting times equal to zero or two with probabilities \(1/2\)
each have mean one and satisfy \(N(t)/t\to1\) almost surely.
Waiting times with \(\mathbb P(X=2^j)=2^{-j}\), \(j\geq1\),
have infinite mean; their counts remain finite at finite times and
satisfy \(N(t)/t\to0\) almost surely.
\item If \(\nu(\{0\})=7/10\) and \(\nu(\{2\})=3/10\), and
\(K_n=\sum_{i=1}^n\mathbf1_{\{X_i=0\}}\), then
\(D_n=|K_n/n-7/10|\) almost surely. At grid size ten,
\(x_{j,10}=0\) for \(1\leq j\leq7\) and \(x_{j,10}=2\)
for \(8\leq j\leq9\); the grid repetitions come from the atoms.
\item On \(\mathcal A=\{0,1,\bot\}\), where \(\bot\) is a third
alphabet symbol, let
\(p(0)=3/4\), \(p(1)=1/4\) and \(p(\bot)=0\). Then
\[
 H=\log4-\tfrac34\log3,\qquad
 \sigma_I^2=\tfrac3{16}(\log3)^2.
\]
A word with \(k\) ones and no forbidden symbol has probability
\(3^{n-k}4^{-n}\) and is typical exactly when
\(|k/n-1/4|\leq\varepsilon/\log3\). Thus its typical-set cardinality is
\(\sum_{0\leq k\leq n,\ |k/n-1/4|\leq\varepsilon/\log3}\binom nk\).
At \(n=4\) and \(\varepsilon=(\log3)/8\), this set has four words
and probability \(27/64\).
\item For a uniform law on \(m\) supported symbols, \(H=\log m\)
and \(\sigma_I^2=0\); every supported word is typical and the
typical set has exactly \(m^n\) words, including when \(m=1\).
\end{enumerate}
\end{hbtheorem}

\begin{hbtheorem}{thm:p15-strong-law-application-boundaries}{proved=annex:p15.strong-law-application-boundaries}{Limits of the extensions and uniformity claims}
\begin{enumerate}
\item Let the i.i.d. symmetric variables assign mass \(2^{-j-1}\)
to each of \(2^j\) and \(-2^j\), for \(j\geq1\). Then
both part expectations are infinite. Almost surely the averages have
no finite or infinite limit; symmetry does not assign them mean zero.
\item All-zero waiting times give \(N(t)=\infty\) at every
nonnegative time. Allowing negative waiting times without a replacement
model also destroys the count interpretation: \(X_n=-1\) gives
\(T_n=-n\) and \(N(t)=\infty\) for every \(t\geq0\).
\item Repeating a single fair \(\{0,1\}\)-valued observation at
every index gives identically distributed dependent data with
\(D_n=1/2\) at every outcome. Separately, \(G_n(x)=\mathbf1_{\{x\geq n\}}\)
are CDFs converging pointwise to the bounded continuous zero function,
but their uniform distance from that function is one. The zero function
is not a CDF.
\item For uniform observations on \([0,1]\), the supremum over
all Borel sets of \(|\nu_n(B)-\nu(B)|\) is one at every outcome
whose observations lie in \([0,1]\). Uniform CDF convergence concerns
half-lines and does not imply this stronger assertion for all Borel sets.
\item If a single bit \(B\), with \(p(0)=3/4\) and \(p(1)=1/4\),
is repeated as \(X_n=B\), then for \(n\geq2\) the marginal product
\(\prod_{i=1}^np(X_i)\) is not the joint probability of the observed
word. Its normalized negative logarithm is the random constant
\(-\log p(B)\), not \(H\). With \(\varepsilon=(\log3)/8\),
the product-law typical set has probability zero under this dependent
process for every \(n\), while two constant words already carry all
of its true probability.
\end{enumerate}
\end{hbtheorem}

\begin{proofboundary}{def:p15-strong-law-applications-next}{Rates and infinite-mean normalization}
The extended-mean law and the complete renewal, empirical-CDF and
finite-alphabet information applications are proved here. The source's
later rates remain separate: Theorem 2.5.11 for finite variance and
Theorem 2.5.12 for the Marcinkiewicz--Zygmund regime, followed by
Feller's infinite-mean normalization Theorem 2.5.13
\cite[pp.~87--90]{DurrettPTEAuthor2019}. The typical-set Chebyshev
bound above is a finite-alphabet calculation and does not discharge
those general rate theorems. The renewal-reward and further applications
in the Section 2.4 exercises also have their own proof obligations.
\end{proofboundary}
